% ============================================================
% PHẦN 2. DỮ LIỆU VÀ PHƯƠNG PHÁP
% ============================================================
\setroadmap{-1}
\section{Dữ liệu và phương pháp}

% ------------------------------------------------------------
% F2.1 Tổng quan pipeline
% ------------------------------------------------------------
\begin{frame}
\frametitle{LegalSCM-RAG: hai pha, năm bước}
\vspace{-0.05cm}
\centering
\begin{tikzpicture}[font=\scriptsize]
  % Hàng 1: ngoại tuyến
  \node[sn=stepData] (raw) at (1.55,0) {\textbf{BLHS}\\406 điều luật};
  \node[sn=stepData,text width=2.95cm] (pre) at (5.15,0) {\textbf{Tiền xử lý}\\(a1) trích xuất rule\\(a2) chuẩn hóa sự kiện};
  \node[sn=stepGraph] (g) at (8.75,0) {\textbf{Bước 1 $\cdot$ Đồ thị}\\nhân quả pháp lý\\5.134 nút};
  \node[sn=stepMemory] (mem) at (12.35,0) {\textbf{Bước 2 $\cdot$ Bộ nhớ}\\nhân quả\\4.698 vector};
  \draw[flow] (raw) -- (pre);
  \draw[flow] (pre) -- (g);
  \draw[flow] (g) -- (mem);
  % Hàng 2: theo từng câu hỏi
  \node[stepnode=schemaGray,text width=1.15cm,minimum height=1.0cm] (q) at (0.7,-2.25) {\textbf{Câu hỏi}\\$q$};
  \node[sn=stepRetrieve,text width=2.2cm] (s3) at (3.1,-2.25) {\textbf{Bước 3}\\Truy xuất đường dẫn\\$\rightarrow P^{*}$};
  \node[sn=stepVerify,text width=2.3cm] (s4) at (5.95,-2.25) {\textbf{Bước 4}\\Xác minh LegalSCM\\theo loại câu hỏi};
  \node[draw=stepVerify,fill=stepVerify!10,shape=diamond,aspect=1.7,inner sep=0.5pt,
        line width=0.7pt,font=\scriptsize\bfseries] (gate) at (8.5,-2.25) {Commit?};
  \node[sn=stepAnswer,text width=2.15cm] (s5) at (13.05,-2.25) {\textbf{Bước 5}\\Sinh câu trả lời\\+ trích dẫn $(a,C)$};
  \draw[flow] (q) -- (s3);
  \draw[flow] (s3) -- (s4);
  \draw[flow] (s4) -- (gate);
  \draw[flow] (gate.north east) to[bend left=10] node[above,align=center,font=\scriptsize]{đạt:\\quyết định SCM} (s5.west);
  \draw[flow,dashed] (gate.south east) to[bend right=10] node[below,align=center,font=\scriptsize]{không đạt:\\ngữ cảnh Bước 3} (s5.west);
  % Nền của hai pha
  \begin{scope}[on background layer]
    \node[fill=schemaLight,rounded corners=4pt,inner sep=5pt,fit=(raw)(mem)] (off) {};
    \node[fill=schemaLight,rounded corners=4pt,inner sep=5pt,fit=(q)(s5)(gate)] (on) {};
  \end{scope}
  \node[anchor=south west,font=\scriptsize\bfseries,text=schemaNavy,inner sep=1pt] at (off.north west) {Ngoại tuyến (một lần)};
  \node[anchor=south west,font=\scriptsize\bfseries,text=schemaNavy,inner sep=1pt] at (on.north west) {Theo từng câu hỏi};
  \draw[flow,dashed,draw=schemaNavy] (off.south) -- node[right,font=\scriptsize,text=schemaNavy]{tài nguyên dùng chung: đồ thị, bộ nhớ, tập rule $\mathcal R$} ($(on.north)+(0,0.42)$);
\end{tikzpicture}

\vspace{0.1cm}
\raggedright
\footnotesize $q\rightarrow P^{*}\rightarrow(z,\mathcal E^{*})\rightarrow(a,C)$:\quad
chọn đường suy luận $\rightarrow$ ra quyết định có xác minh $\rightarrow$ sinh câu trả lời kèm căn cứ.

\vspace{0.1cm}
\keymsg{Điểm mới: bước xác minh LegalSCM (Legal Structural Causal Model, mô hình quy tắc quy phạm ba giá trị)
được đặt có chọn lọc giữa truy xuất đường dẫn và sinh câu trả lời.}
\end{frame}

% ------------------------------------------------------------
% F2.2 Dữ liệu nguồn
% ------------------------------------------------------------
\setroadmap{0}
\begin{frame}
\frametitle{Dữ liệu nguồn: 406 điều của Bộ luật Hình sự}
\small
\begin{itemize}
  \item \textbf{Văn bản:} VBHN số 135/VBHN-VPQH -- Bộ luật Hình sự (VPQH, 05/09/2025).
  \item \textbf{Dữ liệu thô:} 406 điều, mã 1--426 (không có 20 mã: Điều 90--107, 292, 317);
        mỗi bản ghi gồm mã điều, tiêu đề và nội dung.
  \item \textbf{Sau trích xuất:} 2.884 rule từ 395 điều; 11 điều không sinh rule nào.
  \item Trung bình 7,30 rule/điều; nhiều nhất 41 rule (Điều 251); 41 chủ thể pháp lý.
  \item Nhãn chất lượng của rule cuối: \code{CURATED\_FROM\_STATUTORY\_CLAUSE} 2.654;
        \code{CURATED} 230.
\end{itemize}

\vspace{0.15cm}
\begin{center}
\kpi[stepData]{406}{điều luật\\(dữ liệu thô)}\hfill
\kpi[stepData]{2.884}{quy tắc\\điều kiện--hệ quả}\hfill
\kpi[stepGraph]{1.814}{sự kiện\\chuẩn hóa}\hfill
\kpi[stepData]{41}{chủ thể\\pháp lý}
\end{center}

\vspace{0.05cm}
\keymsg{Đơn vị tri thức là rule điều kiện--hệ quả gắn với điều luật nguồn, không phải đoạn văn.}
\end{frame}

% ------------------------------------------------------------
% F2.3a Trích xuất quy tắc
% ------------------------------------------------------------
\begin{frame}
\frametitle{Tiền xử lý (a1): trích xuất quy tắc từ điều luật bằng LLM}
\begin{columns}[T,onlytextwidth]
\column{0.46\textwidth}
\footnotesize
\begin{itemize}
  \item Mỗi điều luật $\rightarrow$ Qwen3-8B (Ollama, temperature 0, đầu ra theo JSON schema).
  \item Mỗi rule gồm: chủ thể, điều kiện, hệ quả và câu luật \code{rule\_text}.
  \item \code{rule\_text} theo mẫu ``Nếu $\langle$điều kiện$\rangle$ thì $\langle$hệ quả$\rangle$.'';
        chỉ dùng nội dung có trong điều luật.
  \item Khử trùng lặp; thử lại tối đa 3 lần khi đầu ra lỗi.
  \item Bộ rule cuối được rà soát trước khi dựng đồ thị.
\end{itemize}

\column{0.50\textwidth}
\centering
\begin{tikzpicture}
  \node[quote=schemaGray,text width=6.3cm] (src) at (0,0)
    {\textbf{Điều 22, khoản 2:}
     \textit{``Vượt quá giới hạn phòng vệ chính đáng là hành vi chống trả rõ ràng quá mức
     cần thiết, không phù hợp với tính chất và mức độ nguy hiểm cho xã hội của hành vi xâm hại.''}};
  \node[stepnode=stepData,text width=3.6cm,minimum height=0.5cm,below=0.3cm of src] (llm)
    {Qwen3-8B + JSON schema};
  \node[quote=stepGraph,text width=6.3cm,below=0.3cm of llm] (rule)
    {\textbf{rule 47} (bản cuối, đã rà soát) -- chủ thể: \textit{Người phòng vệ}\\
     \textbf{Điều kiện:} chống trả rõ ràng quá mức cần thiết và không phù hợp với tính chất,
     mức độ nguy hiểm của hành vi xâm hại\\
     \textbf{Hệ quả:} được xác định là vượt quá giới hạn phòng vệ chính đáng};
  \draw[flow] (src) -- (llm);
  \draw[flow] (llm) -- (rule);
\end{tikzpicture}
\end{columns}

\vspace{0.12cm}
\keymsg{Mỗi rule là một cặp điều kiện $\rightarrow$ hệ quả gắn với điều luật nguồn.}
\end{frame}

% ------------------------------------------------------------
% F2.3b Chuẩn hóa sự kiện
% ------------------------------------------------------------
\begin{frame}
\frametitle{Tiền xử lý (a2): chuẩn hóa sự kiện để nối các điều luật}
\begin{columns}[T,onlytextwidth]
\column{0.57\textwidth}
\footnotesize
\begin{itemize}
  \item Nhiều cách diễn đạt $\rightarrow$ một \code{event\_id}: ``phải chịu TNHS'', ``chịu TNHS'',
        ``bị truy cứu TNHS'' $\rightarrow$ \textit{Chịu trách nhiệm hình sự}.
  \item Không gộp khái niệm chỉ gần nghĩa: ``Miễn TNHS'' $\neq$ ``Không phải chịu TNHS''.
  \item Điểm ứng viên
        $s_{\text{lex}}=0{,}45\,\text{Jaccard}+0{,}35\,\text{SeqMatch}+0{,}20\,\text{containment}$\\
        \dien độ trùng từ + độ giống chuỗi + mức bao hàm.
  \item LLM chọn \code{USE\_EXISTING} hoặc \code{CREATE\_NEW}; chỉ dùng sự kiện có sẵn khi
        điểm $\ge 0{,}60$ và đúng vai trò.
  \item Kết quả: 1.814 sự kiện (1.416 ở phía điều kiện, 412 ở phía hệ quả, 14 ở cả hai).
\end{itemize}

\column{0.39\textwidth}
\centering
\begin{tikzpicture}
  \node[quote=schemaGray,text width=2.25cm] (e47) at (0,0.95)
    {\textbf{rule 47 $\cdot$ hệ quả}\\``được xác định là vượt quá giới hạn PVCĐ''};
  \node[quote=schemaGray,text width=2.25cm] (c106) at (0,-0.95)
    {\textbf{rule 106 $\cdot$ điều kiện}\\``người phạm tội phạm tội trong trường hợp vượt quá giới hạn PVCĐ''};
  \node[bridgenode=stepGraph,text width=1.85cm] (ev) at (3.15,0)
    {\textbf{Vượt quá giới hạn phòng vệ chính đáng}};
  \draw[flow] (e47.east) -- (ev.west);
  \draw[flow] (c106.east) -- (ev.west);
  \node[tag,below=2pt of ev,align=center] {cùng một sự kiện\\(cầu nối)};
\end{tikzpicture}
\end{columns}

\vspace{0.1cm}
\keymsg{Chuẩn hóa nối Điều 22 với Điều 51 qua một sự kiện chung.}
\end{frame}

% ------------------------------------------------------------
% F2.4 Bước 1: đồ thị
% ------------------------------------------------------------
\setroadmap{1}
\begin{frame}
\frametitle{Bước 1: đồ thị nhân quả pháp lý dị thể}
\begin{columns}[T,onlytextwidth]
\column{0.45\textwidth}
\centering
\begin{tikzpicture}[font=\scriptsize]
  \node[tnode=schemaGray] (art) at (0,0) {ARTICLE $a_j$};
  \node[tnode=stepGraph] (rule) at (0,-1.0) {RULE $r_j$};
  \node[tnode=schemaGray] (subj) at (2.75,0) {SUBJECT};
  \node[tnode=stepGraph] (u) at (-1.6,-1.95) {EVENT $u_j$};
  \node[tnode=stepGraph] (v) at (1.6,-1.95) {EVENT $v_j$};
  \draw[flow] (art) -- node[left,font=\scriptsize]{HAS\_RULE} (rule);
  \draw[flow] (rule.east) -- node[below right=-2pt,font=\scriptsize]{HAS\_SUBJECT} (subj.south west);
  \draw[flow] (rule) -- node[left=2pt,font=\scriptsize]{HAS\_CONDITION} (u);
  \draw[flow] (rule) -- node[right=2pt,font=\scriptsize]{HAS\_EFFECT} (v);
  \draw[causes=stepGraph,line width=1.3pt] (u) -- node[below,font=\scriptsize\bfseries,text=stepGraph!80!black]{CAUSES} (v);
\end{tikzpicture}

\vspace{0.05cm}
\raggedright\footnotesize
$r_j:\;u_j\xrightarrow{\;r_j,\,a_j\;}v_j$\quad
\dien mỗi rule là một cạnh CAUSES mang mã rule và điều luật nguồn.

\vspace{0.08cm}
$A\xrightarrow{r_1}B,\ B\xrightarrow{r_2}C\ \Rightarrow\ A\rightarrow B\rightarrow C$\quad
\dien hệ quả rule trước = điều kiện rule sau $\Rightarrow$ chuỗi.

\column{0.51\textwidth}
\footnotesize
\renewcommand{\arraystretch}{1.02}
\begin{tabular}{@{}l r@{}}
\toprule
Tổng số nút & \textbf{5.134}\\
\quad ARTICLE $\cdot$ RULE & 395 $\cdot$ 2.884\\
\quad EVENT $\cdot$ SUBJECT & 1.814 $\cdot$ 41\\
Cạnh (5 loại $\times$ 2.884) & \textbf{14.420}\\
Cặp sự kiện CAUSES duy nhất & 2.122\\
Sự kiện cầu nối $\cdot$ chuỗi hai bước & 14 $\cdot$ 112\\
Đồ thị sự kiện & DAG\\
Thành phần liên thông yếu & 248\\
\bottomrule
\end{tabular}

\vspace{0.1cm}
\begin{itemize}
  \item NetworkX MultiDiGraph: cạnh CAUSES giữ mã rule và mã điều luật
        $\Rightarrow$ nhiều rule cùng hỗ trợ một cặp sự kiện mà không mất nguồn gốc.
\end{itemize}
\end{columns}

\vspace{0.1cm}
\keymsg{Đồ thị phục vụ truy xuất và phân tích tô-pô; ngữ nghĩa suy luận do tập rule $\mathcal R$ quyết định.}
\end{frame}

% ------------------------------------------------------------
% F2.5 Lát cắt đồ thị quanh ví dụ
% ------------------------------------------------------------
\begin{frame}
\frametitle{Lát cắt thật của đồ thị quanh ví dụ xuyên suốt}
\vspace{-0.1cm}
\centering
\begin{tikzpicture}[font=\scriptsize,
  lbl/.style={font=\scriptsize,fill=white,inner sep=0.8pt,align=center,text=stepGraph!70!black}]
  % Điều kiện
  \node[evnode=stepGraph,text width=2.8cm] (c1) at (1.5,0) {Chống trả cần thiết để bảo vệ lợi ích chính đáng};
  \node[evnode=stepGraph,text width=2.8cm] (c2) at (1.5,-1.2) {Gây thiệt hại nhỏ hơn để tránh thiệt hại lớn hơn};
  \node[evnode=stepGraph,text width=2.8cm,line width=1.1pt] (c3) at (1.5,-2.55) {Chống trả rõ ràng quá mức cần thiết};
  % Cầu nối
  \node[bridgenode=stepGraph,text width=2.6cm] (b1) at (7.0,0) {Phòng vệ chính đáng};
  \node[bridgenode=stepGraph,text width=2.6cm] (b2) at (7.0,-1.2) {Tình thế cấp thiết};
  \node[bridgenode=stepGraph,text width=2.6cm,line width=1.1pt] (b3) at (7.0,-2.55) {Vượt quá giới hạn phòng vệ chính đáng};
  % Hệ quả
  \node[evnode=stepGraph,text width=2.9cm] (e1) at (12.55,-0.6) {Không bị coi là tội phạm\\{\color{schemaGray}$\leftarrow$ 5 rule (Điều 8, 22--25)}};
  \node[evnode=stepGraph,text width=2.9cm] (e2) at (12.55,-1.95) {Chịu trách nhiệm hình sự\\{\color{schemaGray}$\leftarrow$ 18 rule}};
  \node[evnode=stepGraph,text width=2.9cm,line width=1.1pt] (e3) at (12.55,-3.15) {Có tình tiết giảm nhẹ TNHS\\{\color{schemaGray}$\leftarrow$ 28 rule (Điều 51, 84)}};
  % Cạnh
  \draw[causes=stepGraph!60,line width=0.6pt] (c1) -- node[lbl,above=1pt]{rule 45 $\cdot$ Điều 22} (b1);
  \draw[causes=stepGraph!60,line width=0.6pt] (c2) -- node[lbl,above=1pt]{rule 49 $\cdot$ Điều 23} (b2);
  \draw[causes=stepGraph,line width=1.4pt] (c3) -- node[lbl,above=1pt]{rule 47 $\cdot$ Điều 22} (b3);
  \draw[causes=stepGraph!60,line width=0.6pt] (b1.east) -- node[lbl,above=1pt,pos=0.45]{rule 46 $\cdot$ Điều 22} (e1.west);
  \draw[causes=stepGraph!60,line width=0.6pt] (b2.east) -- node[lbl,below=1pt,pos=0.45]{rule 50 $\cdot$ Điều 23} (e1.west);
  \draw[causes=stepGraph!60,line width=0.6pt] (b3.east) -- node[lbl,above=1pt,pos=0.5]{rule 48 $\cdot$ Điều 22} (e2.west);
  \draw[causes=stepGraph,line width=1.4pt] (b3.east) -- node[lbl,below=1pt,pos=0.5]{rule 106 $\cdot$ Điều 51} (e3.west);
  \node[tag,below=1pt of c3] {\textbf{đường của ví dụ} \exid};
\end{tikzpicture}

\vspace{0.05cm}
\raggedright
\footnotesize
\begin{itemize}
  \item 3 trong 14 sự kiện cầu nối nằm trong lát cắt này; có \textbf{rẽ nhánh} (một cầu nối $\rightarrow$
        hai hệ quả) và \textbf{hội tụ} (hai đường $\rightarrow$ \textit{Không bị coi là tội phạm}, tức tuyển giữa các rule).
\end{itemize}

\vspace{0.05cm}
\keymsg{Sự kiện cầu nối biến hai điều luật rời rạc thành một chuỗi suy luận hai bước.}
\end{frame}

% ------------------------------------------------------------
% F2.6 Bước 2: bộ nhớ nhân quả
% ------------------------------------------------------------
\setroadmap{2}
\begin{frame}
\frametitle{Bước 2: bộ nhớ nhân quả cho tìm kiếm ngữ nghĩa}
\centering
\begin{tikzpicture}[font=\scriptsize]
  \node[stepnode=stepMemory,text width=4.6cm,minimum height=0.72cm,align=left] (rm) at (0,0.43)
    {\textbf{RULE memory}: chủ thể, điều kiện, hệ quả, điều luật};
  \node[stepnode=stepMemory,text width=4.6cm,minimum height=0.72cm,align=left] (em) at (0,-0.43)
    {\textbf{EVENT memory}: tên, alias, vai trò, ngữ cảnh};
  \node[stepnode=stepMemory,text width=1.75cm] (enc) at (4.25,0) {\textbf{BGE-M3}\\1.024 chiều};
  \node[stepnode=stepMemory,text width=1.5cm] (l2) at (6.65,0) {Chuẩn hóa\\L2};
  \node[stepnode=stepMemory,text width=1.9cm] (idx) at (9.15,0) {\textbf{FAISS}\\IndexFlatIP};
  \draw[flow] (rm.east) -- (enc);
  \draw[flow] (em.east) -- (enc);
  \draw[flow] (enc) -- (l2);
  \draw[flow] (l2) -- (idx);
  \node[tag,anchor=west] at (10.35,0) {4.698 vector};
\end{tikzpicture}

\vspace{0.1cm}
\raggedright
\begin{columns}[T,onlytextwidth]
\column{0.49\textwidth}
\footnotesize
\textbf{Trực giác:} mỗi rule và mỗi sự kiện thành một vector; câu hỏi tìm các vector gần nhất.

\vspace{0.1cm}
\centering
$\hat{\mathbf z}_j=\dfrac{\mathbf z_j}{\|\mathbf z_j\|_2},\qquad
s(q,c_j)=\hat{\mathbf q}^{\top}\hat{\mathbf z}_j=\cos(\mathbf q,\mathbf z_j)$

\vspace{0.05cm}
\raggedright
\dien sau chuẩn hóa, tích vô hướng chính là cosine; IndexFlatIP tìm kiếm chính xác, không xấp xỉ.

\vspace{0.05cm}
{\scriptsize\color{schemaGray} BGE-M3: Chen và cs. (2024).}

\column{0.47\textwidth}
\footnotesize
\begin{itemize}
  \item 4.698 vector = 2.884 RULE + 1.814 EVENT.
  \item Vai trò EVENT: CONDITION 1.402, EFFECT 398, BRIDGE 14.
  \item \code{memory\_id} trùng vị trí vector; kiểm tra nhất quán CSV -- embedding -- FAISS.
\end{itemize}
\vspace{0.05cm}
\exbox{stepMemory}{0.92\linewidth}{ba EVENT xếp đầu \#2956 (CONDITION), \#2957 (BRIDGE),
\#3047 (EFFECT) là đúng ba sự kiện của chuỗi; RULE 106 (Điều 51) gần câu hỏi nhất, cos = 0,653.}
\end{columns}
\end{frame}

% ------------------------------------------------------------
% F2.7 Bước 3: truy xuất đường dẫn
% ------------------------------------------------------------
\setroadmap{3}
\begin{frame}
\frametitle{Bước 3: truy xuất đường dẫn theo hướng câu hỏi}
\begin{columns}[T,onlytextwidth]
\column{0.47\textwidth}
\small
\begin{enumerate}
  \item Mã hóa câu hỏi bằng BGE-M3.
  \item FAISS lấy 100 RULE/EVENT gần nhất.
  \item Chọn seed: EVENT top-8 ($s\ge0{,}20$) và RULE top-8 ($s\ge0{,}15$).
  \item Suy hướng từ ý định câu hỏi: xuôi, ngược hoặc cả hai.
  \item BFS tối đa 2 hop, không lặp nút, tối đa 30 đường mỗi seed.
  \item Gom rule từ seed, đường và kết quả trực tiếp $\rightarrow$ xếp hạng $\rightarrow$ 12 bằng chứng.
\end{enumerate}

\column{0.49\textwidth}
\centering
\exbox{stepRetrieve}{0.92\linewidth}{Cụm ``bắt nguồn từ'' $\Rightarrow$ hướng \textbf{ngược} (backward).}

\vspace{0.15cm}
\begin{tikzpicture}
  \node[funnel=10,minimum width=6.4cm] (f1) at (0,0) {100 kết quả ngữ nghĩa (RULE + EVENT)};
  \node[funnel=20,minimum width=5.7cm] (f2) at (0,-0.6) {8 seed sự kiện + 8 rule trực tiếp};
  \node[funnel=30,minimum width=5.0cm] (f3) at (0,-1.2) {33 đường ứng viên, 1 đường đủ 2 bước};
  \node[funnel=42,minimum width=4.0cm] (f4) at (0,-1.8) {36 rule ứng viên};
  \node[funnel=55,minimum width=3.2cm,font=\scriptsize\bfseries] (f5) at (0,-2.4) {12 bằng chứng cuối};
  \foreach \a/\b in {f1/f2,f2/f3,f3/f4,f4/f5} \draw[flow,shorten >=0pt,shorten <=0pt] (\a) -- (\b);
\end{tikzpicture}
\end{columns}

\vspace{0.25cm}
\keymsg{Khác RAG thường: ngữ cảnh gồm sự kiện, rule và quan hệ có hướng nối chúng, không chỉ các đoạn văn độc lập.}
\end{frame}

% ------------------------------------------------------------
% F2.8 Xếp hạng đường dẫn
% ------------------------------------------------------------
\begin{frame}
\frametitle{Xếp hạng đường dẫn: ưu tiên đường đủ bước qua cầu nối}
\footnotesize
\textbf{Trực giác:} đường tốt bắt đầu từ seed gần câu hỏi, được rule hỗ trợ, đi qua sự kiện cầu nối
và đủ số bước.

\vspace{0.1cm}
\begin{columns}[T,onlytextwidth]
\column{0.55\textwidth}
\footnotesize
\textbf{Bài báo:} $P^{*}=\arg\max_{P_i\in\mathcal P(q)}S(P_i,q)$\\
\dien chọn đường có điểm phù hợp ngữ nghĩa và cấu trúc cao nhất.

\vspace{0.12cm}
\textbf{Điểm được cài đặt:}\\[2pt]
$S_{\text{path}}=\operatorname{clip}_{[0,1]}\big(0{,}54\,s_{\text{seed}}+0{,}28\,s_{\text{sup}}
+b_{\text{bridge}}+b_{\text{hop}}\big)$\\[3pt]
$b_{\text{bridge}}=\min(0{,}18;\,0{,}06\,N_{\text{bridge}}),\quad b_{\text{hop}}=0{,}18$ nếu đủ 2 bước\\[2pt]
\dien độ gần của seed chiếm phần lớn; hai khoản thưởng ưu tiên đường hai bước qua cầu nối.

\column{0.41\textwidth}
\exbox{stepRetrieve}{0.92\linewidth}{%
$S(P_0)=0{,}54\cdot0{,}602+0{,}28\cdot0{,}5+0{,}06+0{,}18=\mathbf{0{,}705}$
$\Rightarrow$ hạng 1/33.\\[2pt]
Đường 1 bước tốt nhất: $S(P_1)=0{,}481$.}

\vspace{0.1cm}
\footnotesize
Điểm bằng chứng của rule kết hợp 5 tín hiệu (Phụ lục A1):
rule 106 = 0,719 (hạng 1), rule 47 = 0,642 (hạng 2).
\end{columns}

\vspace{0.15cm}
\keymsg{Với ví dụ, đường đúng đứng đầu; trên toàn benchmark, xếp hạng đường vẫn là nút thắt (Phần 4).}
\end{frame}

% ------------------------------------------------------------
% F2.9 Ngữ nghĩa ba giá trị
% ------------------------------------------------------------
\setroadmap{4}
\begin{frame}
\frametitle{Bước 4: mỗi sự kiện pháp lý có ba trạng thái}
\begin{columns}[T,onlytextwidth]
\column{0.50\textwidth}
\footnotesize
\begin{itemize}
  \item \textbf{Trực giác:} thông tin pháp lý thường chưa đầy đủ -- ``chưa biết'' khác ``sai''.
  \item $X_i\in\{1,0,\bot\}$ = TRUE / FALSE / UNKNOWN; $\bot$ là thiếu hoặc xung đột hỗ trợ,
        không phải giá trị nằm giữa 0 và 1.
  \item Rule kích hoạt khi \textbf{mọi} điều kiện thỏa (AND trong rule); hệ quả đúng khi
        \textbf{ít nhất một} rule kích hoạt (OR giữa các rule); lặp đến điểm bất động.
  \item Quy ước cài đặt: sự kiện gốc chưa cho $\rightarrow$ UNKNOWN; hệ quả không có rule nào
        kích hoạt $\rightarrow$ FALSE (thế giới đóng); xung đột $\rightarrow$ UNKNOWN.
\end{itemize}
{\scriptsize\color{schemaGray} Nền tảng: Kleene (1952); Fitting (1985).}

\column{0.46\textwidth}
\exbox{stepVerify}{0.92\linewidth}{điểm bất động factual (--: chưa suy ra).}

\vspace{0.1cm}
\scriptsize
\setlength{\tabcolsep}{3pt}
\renewcommand{\arraystretch}{1.25}
\newcommand{\tv}{\cellcolor{schemaGreen!22}TRUE}
\newcommand{\un}{\cellcolor{schemaGray!12}--}
\begin{tabular}{@{}>{\raggedright\arraybackslash}p{1.55cm}
                >{\centering\arraybackslash}p{1.45cm}
                >{\centering\arraybackslash}p{1.45cm}
                >{\centering\arraybackslash}p{1.45cm}@{}}
\toprule
\textbf{Vòng} & \textbf{Chống trả rõ ràng quá mức} & \textbf{Vượt quá giới hạn PVCĐ} & \textbf{Có tình tiết giảm nhẹ TNHS}\\
\midrule
Ngữ cảnh & \tv & \un & \un\\
Sau rule 47 & \tv & \tv & \un\\
Sau rule 106 & \tv & \tv & \tv\\
\bottomrule
\end{tabular}

\vspace{0.08cm}
Hội tụ sau 3 vòng lặp; mọi hop có rule kích hoạt (47, 106); không có xung đột
$\Rightarrow$ outcome TRUE.
\end{columns}

\vspace{0.05cm}
\keymsg{Kết luận chỉ được chấp nhận khi chuỗi rule thực sự suy ra được hệ quả.}
\end{frame}

% ------------------------------------------------------------
% F2.10 Biểu diễn thưa
% ------------------------------------------------------------
\begin{frame}
\frametitle{Biểu diễn thưa: AND trong một rule, OR giữa các rule}
\begin{columns}[T,onlytextwidth]
\column{0.55\textwidth}
\footnotesize
$(t_i,f_i)=(1,0),(0,1),(0,0)$ khi $X_i=1,0,\bot$\quad \dien hai kênh hỗ trợ đúng/sai.

\vspace{0.14cm}
$\mathbf b=\mathbb I[\mathbf B^{+}\mathbf t+\mathbf B^{-}\mathbf f=\mathbf k],\qquad
 \mathbf a=\mathbf b\odot(\mathbf 1-\mathbf e)$\\[1pt]
\dien rule kích hoạt khi đủ mọi literal trong thân ($\mathbf k$ = số literal) và không bị
ngoại lệ $\mathbf e$ chặn.

\vspace{0.14cm}
$\mathbf d^{+}=\mathbb I[\mathbf H^{+}\mathbf a\ge1],\ \ \mathbf d^{-}=\mathbb I[\mathbf H^{-}\mathbf a\ge1]$\quad
\dien có hỗ trợ khi $\ge1$ rule kích hoạt.

\vspace{0.14cm}
$\mathbf x^{F}=\operatorname{Fix}(\Phi_{\mathcal R},\mathbf x^{(0)})$\quad
\dien lặp đến khi trạng thái không đổi.

\vspace{0.14cm}
{\scriptsize\color{schemaGray} Nền tảng: Sato và cs. (2020); Sakama và cs. (2021).}

\column{0.41\textwidth}
\begin{block}{Hội và tuyển}
\footnotesize
$(A\wedge B)\to Y$: một rule có hai literal.\\
$(A\vee B)\to Y$: tách thành $A\to Y$ và $B\to Y$.
\end{block}
\vspace{0.05cm}
\begin{block}{Quy mô trong thực nghiệm}
\footnotesize
$\mathbf B^{+}\in\{0,1\}^{2884\times1814}$, $\mathbf H^{+}\in\{0,1\}^{1814\times2884}$:
mỗi ma trận 2.884 phần tử khác 0; $\mathbf B^{-}$, $\mathbf H^{-}$, ngoại lệ rỗng.
187 sự kiện có $\ge2$ rule hỗ trợ (tuyển).
\end{block}
\end{columns}

\vspace{0.15cm}
{\footnotesize\textcolor{schemaRed}{\textbf{Giới hạn:}} mỗi rule hiện chỉ có một tiền đề dương
(0 rule hội) $\Rightarrow$ phép hội chưa được kiểm chứng thực nghiệm.}
\end{frame}

% ------------------------------------------------------------
% F2.11 Ba route xác minh
% ------------------------------------------------------------
\begin{frame}
\frametitle{Ba route xác minh theo loại câu hỏi}
\vspace{-0.05cm}
\footnotesize
\renewcommand{\arraystretch}{1.12}
\setlength{\tabcolsep}{4pt}
\begin{tabular}{@{}>{\raggedright\arraybackslash}p{0.17\textwidth}
                >{\raggedright\arraybackslash}p{0.26\textwidth}
                >{\raggedright\arraybackslash}p{0.42\textwidth}
                >{\centering\arraybackslash}p{0.07\textwidth}@{}}
\toprule
\textbf{Route} & \textbf{Dạng câu hỏi} & \textbf{Cách xác minh} & \textbf{Số câu}\\
\midrule
{\color{stepVerify}\textbf{FACTUAL}} & ``$A$ có dẫn đến $Y$?'' (hỏi xuôi hoặc ngược)
  & Gán seed = TRUE; yêu cầu outcome TRUE, mọi hop có rule kích hoạt, không xung đột. & \textbf{230}\\
{\color{stepVerify}\textbf{DIRECT}} & ``$A$ trực tiếp gây ra $Y$?''
  & Ground hai đầu mút ($\ge0{,}80$), rồi kiểm tra cạnh $A\to Y$. & \textbf{13}\\
{\color{stepVerify}\textbf{INTERVENTION}} & ``Bỏ qua bước $m$, $Y$ còn xảy ra?''
  & Ground sự kiện trung gian ($\ge0{,}80$), chạy $do(X_m=0)$, suy luận lại điểm bất động. & \textbf{7}\\
\bottomrule
\end{tabular}

\vspace{0.2cm}
\begin{itemize}
  \item \textbf{FACTUAL} -- \exid{} (ví dụ xuyên suốt) $\Rightarrow$ \labSup.
  \item \textbf{DIRECT} -- \texttt{BLHS\_MH\_0221}: ``Nhận định \textit{Chống trả rõ ràng quá mức cần thiết}
        trực tiếp gây ra \textit{Chịu trách nhiệm hình sự} mà không cần sự kiện trung gian là đúng hay sai?''\\
        Hai đầu mút được ground với độ tin cậy 1,0; đồ thị chỉ có đường qua
        \textit{Vượt quá giới hạn PVCĐ} (rule 47 $\rightarrow$ rule 48, xem lát cắt đồ thị) $\Rightarrow$ \labRej.
\end{itemize}

\vspace{0.2cm}
{\scriptsize\color{schemaGray} Router dựa trên dấu hiệu từ vựng của câu hỏi; chỉ đường chính (path 0) được xác minh cấu trúc.}
\end{frame}

% ------------------------------------------------------------
% F2.12 do() khác xóa nút
% ------------------------------------------------------------
\begin{frame}
\frametitle{Can thiệp $do(X_m=0)$ không phải là xóa nút}
\vspace{-0.1cm}
\begin{columns}[T,onlytextwidth]
\column{0.48\textwidth}
\centering
{\footnotesize\textbf{Xóa nút (chẩn đoán tô-pô)}}\\[1pt]
\begin{tikzpicture}[font=\scriptsize]
  \node[tnode=schemaGray] (a) at (0,0) {$A$};
  \node[tnode=schemaGray] (m) at (1.7,0.45) {$m$};
  \node[tnode=schemaGray] (d) at (1.7,-0.45) {$D$};
  \node[tnode=schemaGray] (y) at (3.4,0) {$Y$};
  \draw[flow] (a) -- (m); \draw[flow] (m) -- (y);
  \draw[flow,dashed] (a) -- (d); \draw[flow,dashed] (d) -- (y);
  \draw[schemaRed,line width=1.2pt] (m.north west) -- (m.south east) (m.north east) -- (m.south west);
  \node[tag,anchor=west] at (3.75,0) {đường khác?};
\end{tikzpicture}\\[1pt]
\footnotesize
$\rho_{\text{top}}=\mathbb I\big[\,Y\ \text{đạt được trong}\ G^{-m}=G[V\setminus\{m\}]\,\big]$\\
\dien chỉ hỏi còn đường khác hay không (giới hạn 3 hop, 30 đường).

\column{0.48\textwidth}
\centering
{\footnotesize\textbf{Can thiệp cấu trúc}}\\[1pt]
\begin{tikzpicture}[font=\scriptsize]
  \node[tnode=stepVerify] (a) at (0,0) {$A$};
  \node[tnode=stepVerify,fill=schemaRed!12] (m) at (1.7,0.45) {$m$ = FALSE};
  \node[tnode=stepVerify] (d) at (1.7,-0.45) {$D$};
  \node[tnode=stepVerify] (y) at (3.4,0) {$Y$};
  \draw[flow,dotted] (a) -- node[midway,font=\scriptsize\bfseries,text=schemaRed]{$\times$} (m);
  \draw[flow] (m) -- (y);
  \draw[flow,dashed] (a) -- (d); \draw[flow,dashed] (d) -- (y);
  \node[tag,anchor=west] at (3.75,0) {suy ra theo rule?};
\end{tikzpicture}\\[1pt]
\footnotesize
$\mathbf x^{CF(m)}=\operatorname{Fix}\big(\Phi_{\mathcal R^{do(X_m=0)}},\mathbf x^{(0)}[X_m\leftarrow0]\big)$\\
\dien cố định $m$ = FALSE, bỏ cơ chế sinh ra $m$, rồi suy luận lại theo luật.
\end{columns}

\vspace{0.1cm}
\exboxof{stepVerify}{0.97\linewidth}{\texttt{BLHS\_MH\_0223}}{(câu hỏi: có thể bỏ qua bước
\textit{Tình thế cấp thiết} để kết luận trực tiếp không?)\\[1pt]
Factual: \textit{Gây thiệt hại nhỏ hơn\ldots} $\rightarrow$ \textit{Tình thế cấp thiết}
$\rightarrow$ \textit{Không bị coi là tội phạm}: đều TRUE.\\
$do(\text{Tình thế cấp thiết}=0)$: \textit{Không bị coi là tội phạm} TRUE $\rightarrow$ \textbf{FALSE}
(rule 49, 50 không còn kích hoạt) $\Rightarrow$ NECESSARY $\Rightarrow$ \labRej{} (0,723; đã commit).}

\vspace{0.08cm}
\keymsg{$G^{-m}\neq do(X_m=0)$ (Pearl, 2009): xóa nút hỏi ``còn đường khác không?'';
can thiệp hỏi ``theo luật, hệ quả còn được suy ra không?''.}
\end{frame}

% ------------------------------------------------------------
% F2.13 Cổng commit
% ------------------------------------------------------------
\begin{frame}
\frametitle{Cổng commit: chỉ tin bộ xác minh khi đủ điều kiện}
\centering
\begin{tikzpicture}[font=\scriptsize]
  \node[stepnode=stepVerify,text width=2.3cm] (v) at (0,0) {\textbf{Bước 4}\\quyết định + audit};
  \node[draw=stepVerify,fill=stepVerify!10,shape=diamond,aspect=1.7,inner sep=0.5pt,
        line width=0.7pt,font=\scriptsize\bfseries] (g) at (3.25,0) {Đạt cổng?};
  \node[stepnode=stepVerify,text width=3.4cm,minimum height=0.55cm] (pass) at (7.0,0.48) {$z_{\text{SCM}}$ + bằng chứng đã xác minh};
  \node[stepnode=schemaGray,text width=3.4cm,minimum height=0.55cm] (fail) at (7.0,-0.48) {giữ ngữ cảnh Bước 3; LLM tự quyết};
  \node[stepnode=stepAnswer,text width=2.2cm] (llm) at (10.9,0) {\textbf{LLM chung}\\$\rightarrow (a,C)$};
  \draw[flow] (v) -- (g);
  \draw[flow] (g.north east) -- node[above left=-1pt,font=\scriptsize]{đạt} (pass.west);
  \draw[flow,dashed] (g.south east) -- node[below left=-1pt,font=\scriptsize]{không đạt} (fail.west);
  \draw[flow] (pass.east) -- (llm);
  \draw[flow,dashed] (fail.east) -- (llm);
\end{tikzpicture}

\vspace{0.12cm}
\raggedright
\begin{columns}[T,onlytextwidth]
\column{0.55\textwidth}
\footnotesize
$G_{\text{commit}}=\mathbb I\big[g_{\text{ground}}\wedge g_{\text{path}}\wedge(s_{\text{SCM}}\ge\tau)
\wedge g_{\text{consistent}}\big]$, với $\tau=0{,}55$\\
\dien chỉ khi cả bốn điều kiện đạt, quyết định của bộ xác minh mới có thẩm quyền.

\vspace{0.1cm}
$z=z_{\text{SCM}}$ nếu $G_{\text{commit}}=1$, ngược lại $z=z_{\text{LLM}}$;\quad
$\mathcal E^{*}=\mathcal E_{\text{verified}}$ hoặc $\mathcal E_{\text{path}}$.

\vspace{0.1cm}
Khi không đạt: không ép \labUnc, không biến chẩn đoán tô-pô thành quyết định SCM.

\column{0.41\textwidth}
\exbox{stepVerify}{0.92\linewidth}{\\[1pt]
\ok\ đường chính hợp lệ (path 0)\\
\ok\ mọi hop có rule kích hoạt\\
\ok\ 0 xung đột; outcome TRUE\\
\ok\ $s_{\text{SCM}}=0{,}577\ge0{,}55$\\[1pt]
$\Rightarrow$ \texttt{VERIFIER\_COMMITTED}}
\end{columns}

\vspace{0.1cm}
\keymsg{Full LegalSCM là hệ lai có chọn lọc: kết luận hình thức khi đủ căn cứ, quay về ngữ cảnh truy xuất khi không.}
\end{frame}

% ------------------------------------------------------------
% F2.14 Bước 5: sinh câu trả lời
% ------------------------------------------------------------
\setroadmap{5}
\begin{frame}
\frametitle{Bước 5: sinh câu trả lời và vết suy luận đầy đủ}
\begin{columns}[T,onlytextwidth]
\column{0.40\textwidth}
\footnotesize
\[
I_q=(q,P^{*},z,\mathcal E^{*}),\qquad (a,C)=f_\theta(I_q)
\]
\dien LLM diễn đạt quyết định và bằng chứng được cung cấp, kèm trích dẫn.

\column{0.56\textwidth}
\footnotesize
\begin{itemize}
  \item Cùng Qwen3-8B, prompt, JSON schema, parser và seed cho mọi hệ thống; không dùng fallback trích xuất.
  \item Trích dẫn chỉ trong tập điều luật được cung cấp.
  \item Ghi audit: route, grounding, can thiệp, commit, fallback, bằng chứng cuối.
\end{itemize}
\end{columns}

\vspace{0.2cm}
{\footnotesize\textbf{Ví dụ \exid{} qua năm bước}}\\[3pt]
\centering
\begin{tikzpicture}[font=\scriptsize]
  \node[stepnode=stepGraph,text width=2.3cm,minimum height=2.05cm,align=left] (t1) at (0,0)
    {\textbf{1 $\cdot$ Đồ thị}\\chuỗi rule 47 (Điều 22) $\rightarrow$ rule 106 (Điều 51) có sẵn.};
  \node[stepnode=stepMemory,text width=2.3cm,minimum height=2.05cm,align=left] (t2) at (2.95,0)
    {\textbf{2 $\cdot$ Bộ nhớ}\\EVENT \#2956, \#2957, \#3047; RULE 106: cos 0,653.};
  \node[stepnode=stepRetrieve,text width=2.3cm,minimum height=2.05cm,align=left] (t3) at (5.9,0)
    {\textbf{3 $\cdot$ Truy xuất}\\hướng ngược; 33 đường; $P_0$ hạng 1 (0,705).};
  \node[stepnode=stepVerify,text width=2.3cm,minimum height=2.05cm,align=left] (t4) at (8.85,0)
    {\textbf{4 $\cdot$ Xác minh}\\factual; 3 vòng; outcome TRUE; 0,577 $\ge$ 0,55 $\Rightarrow$ commit.};
  \node[stepnode=stepAnswer,text width=2.3cm,minimum height=2.05cm,align=left] (t5) at (11.8,0)
    {\textbf{5 $\cdot$ Trả lời}\\\labSup; ``Sự kiện trước là `Chống trả rõ ràng quá mức cần thiết'.''; Điều 22, 51; 2,83 giây.};
  \foreach \a/\b in {t1/t2,t2/t3,t3/t4,t4/t5} \draw[flow] (\a) -- (\b);
\end{tikzpicture}
\end{frame}
